\documentclass{article}

\usepackage[T1]{fontenc}
\usepackage[utf8]{inputenc}
\usepackage[polish]{babel}
\usepackage{graphicx}
\usepackage[a4paper,left=3cm,right=3cm,top=2.5cm,bottom=2.5cm]{geometry}
\usepackage{comment}
\usepackage{parskip}
\usepackage[utf8]{inputenc}
\usepackage{booktabs}  
\usepackage{makecell}
\usepackage{float}
\usepackage{graphicx}
\usepackage{amsmath}
\usepackage{multirow} 
\usepackage{amssymb}
\usepackage{graphicx}
\usepackage{booktabs}
\usepackage{hyperref}
\usepackage{xcolor}
\usepackage{listings}
\usepackage{array}
\usepackage{geometry}
\geometry{margin=1in}

\lstset{
    basicstyle=\ttfamily\small,
    breaklines=true,
    columns=fullflexible,
    frame=single,
    numbers=left,
    numberstyle=\tiny,
    keywordstyle=\color{blue},
    commentstyle=\color{gray},
    stringstyle=\color{red},
    backgroundcolor=\color{gray!5}
}

\title{Zwiększanie odporności rozpoznawania twarzy na okluzję: \\ modyfikacje IResNet50 z mechanizmami uwagi \\
\large Raport projektowy na zajęcia\\
\large Studio Projektowe}

\author{Eliza Petrycka \\ Jakub Mikołajczyk}
\date{Styczeń 2026}

\begin{document}
\maketitle
\newpage

\section{Streszczenie (Abstract)}
\paragraph{}Systemy rozpoznawania tożsamości na podstawie obrazu twarzy osiągają imponujące wyniki w warunkach kontrolowanych dzięki nowoczesnym głębokim sieciom neuronowym i funkcjom strat opartym na marginesach kątowych, takich jak ArcFace. Ich odporność jednak zazwyczaj spada w obecności okluzji - kluczowego wyzwania w rzeczywistych aplikacjach, w tym systemach monitoringu, biometrii bezpieczeństwa i urządzeniach mobilnych. Niniejsza praca analizuje wpływ częściowej okluzji twarzy, szczególnie zasłonięcia obszaru oczu, na dokładność identyfikacji osób oraz ocenia skuteczność modyfikacji architektur zwiększających odporność na tego typu zakłócenia.

Implementujemy i ewaluujemy trzy warianty architektur opartych na IResNet50: model bazowy, model z rozproszoną uwagą CBAM w każdym bloku rezydualnym oraz model z pomocniczą gałęzią Transformer. Ewaluacja na testowym zbiorze CASIA-WebFace z deterministyczną okluzją 20px wskazuje, że rozproszona uwaga poprawia odporność na zasłonięcie: CBAM\_Block osiąga 90.16\% Rank-1, przewyższając baseline (87.11\%) o 3.05 pp. Wariant Transformer uzyskuje 85.00\% Rank-1.

Analiza potwierdza, że hierarchiczna filtracja zakłóceń na wielu poziomach abstrakcji cech jest skutecznym kierunkiem zwiększania odporności modeli rozpoznawania twarzy w scenariuszach z okluzją. Proponowane modyfikacje oparte na CBAM zachowują prostotę wdrożenia przy jednoczesnej poprawie jakości identyfikacji.


\section{Wstęp}
\paragraph{}Rozpoznawanie tożsamości na podstawie obrazu twarzy (ang. Face Recognition) jest jednym z kluczowych zagadnień współczesnej wizji komputerowej i znajduje szerokie zastosowanie w systemach bezpieczeństwa, monitoringu wizyjnego oraz urządzeniach konsumenckich. W ostatnich latach, dzięki rozwojowi głębokich sieci neuronowych oraz funkcji strat opartych na marginesach kątowych, takich jak ArcFace, nowoczesne modele typu state-of-the-art osiągają bardzo wysoką skuteczność w warunkach kontrolowanych i na standardowych zbiorach testowych.

Jednocześnie skuteczność systemów rozpoznawania twarzy spada w obecności zakłóceń charakterystycznych dla rzeczywistych scenariuszy. Szczególnym wyzwaniem jest częściowa okluzja twarzy, w tym zasłonięcie obszaru oczu spowodowane noszeniem okularów, masek ochronnych lub ograniczoną widocznością w systemach monitoringu. Chociaż złożone obliczeniowo modele typu state-of-the-art wykazują relatywnie wysoką odporność na tego typu zakłócenia, ich lżejsze warianty, projektowane z myślą o zastosowaniach mobilnych, jak również starsze architektury, często charakteryzują się znacznym spadkiem skuteczności identyfikacji w warunkach okluzji.

Celem niniejszej pracy jest analiza wpływu częściowej okluzji twarzy na skuteczność
systemów rozpoznawania tożsamości oraz zaproponowanie modyfikacji architektury
opartej na mechanizmach atencji, zwiększającej odporność modeli na tego typu zakłócenia.

\section{Stan wiedzy}
\subsection{Okluzja jako źródło utraty informacji i szumu w rozpoznawaniu twarzy}

\paragraph{}Fundamentem dla zrozumienia metod radzenia sobie z niepełnymi danymi wizualnymi jest zdefiniowanie natury problemu okluzji. Jak wskazują Zeng i in. [1] w swoim przeglądzie technik rozpoznawania twarzy, okluzja stanowi dwojakie wyzwanie dla systemów biometrycznych. Po pierwsze, powoduje utratę informacji dyskryminacyjnych w obszarach zasłoniętych. Po drugie, i co często groźniejsze, wprowadza szum (błędne cechy) pochodzący z obiektu przysłaniającego (np. maski, dłoni), który może zostać błędnie zinterpretowany przez sieć jako cecha tożsamościowa.

Zeng i in. proponują systematykę podejść badawczych, dzieląc je na trzy główne kategorie: (1) odporną ekstrakcję cech (ang. Occlusion Robust Feature Extraction), która dąży do uzyskania reprezentacji niewrażliwej na zakłócenia, często poprzez skupienie się na cechach lokalnych zamiast globalnych; (2) rozpoznawanie świadome okluzji (ang. Occlusion Aware Face Recognition), gdzie mechanizmy detekcji lub maskowania dynamicznie wykluczają zaszumione obszary z procesu decyzyjnego; oraz (3) rekonstrukcję okluzji (ang. Occlusion Recovery), polegającą na wizualnym odtworzeniu brakujących fragmentów [1].

W kontekście tej klasyfikacji, współczesna literatura coraz częściej odchodzi od prostych metod rekonstrukcji na rzecz zaawansowanego modelowania relacji przestrzennych (podejście typu 1 i 2). Klasyczne sieci splotowe (CNN), choć skuteczne w ekstrakcji cech lokalnych, mogą napotykać trudności w modelowaniu globalnych zależności geometrycznych, gdy część danych wejściowych ulega zmianie
\subsection{Modelowanie relacji przestrzennych w warunkach niepełnych danych}

\paragraph{}Istotnym obszarem badań w dziedzinie rozpoznawania twarzy jest opracowanie metod odpornych na degradację lub zmianę cech biometrycznych, do czego dochodzi zarówno w przypadku okluzji (przysłonięcia), jak i naturalnych procesów starzenia. Współczesna literatura wskazuje, że klasyczne sieci splotowe (CNN) są skuteczne w ekstrakcji cech lokalnych, jednak mogą napotykać trudności w modelowaniu globalnych zależności geometrycznych, zwłaszcza gdy część danych wejściowych ulega degradacji lub zasłonięciu.

W tym kontekście obiecującym kierunkiem jest wykorzystanie sieci typu Transformer do modelowania długozasięgowych zależności przestrzennych (ang. long-range dependencies). Jak wykazują Prakash i Rai [2], relacje przestrzenne pomiędzy kluczowymi elementami twarzy pozostają stosunkowo stałe, nawet gdy powierzchowne cechy (tekstura skóry, zmarszczki, proporcje) ulegają znacznym zmianom w czasie. Autorzy ci proponują hybrydowe podejście, w którym Transformer nie zastępuje głównej sieci ekstrakcyjnej (backbone CNN), lecz wykorzystywany jest jako pomocnicza funkcja straty (ang. auxiliary loss), współtworząc mechanizm Transformer-Enhanced Loss.

\paragraph{}W proponowanej przez nich architekturze, mapa cech z ostatniej warstwy splotowej jest traktowana jako sekwencja wektorów kontekstowych. Mechanizm atencji pozwala sieci skupić się na zachowanych relacjach między widocznymi fragmentami twarzy, co mityguje wpływ zakłóceń. W badaniach na zbiorach o dużej wariancji wiekowej (CA-LFW, AgeDB), połączenie standardowej straty metrycznej (np. ArcFace) ze stratą transformatorową pozwoliło na osiągnięcie wyników typu State-of-the-Art [2].

Dla problematyki okluzji wnioski te są kluczowe: sugerują, że skuteczna identyfikacja twarzy częściowo przysłoniętej nie musi polegać wyłącznie na rekonstrukcji brakujących fragmentów (inpainting), lecz może opierać się na silniejszym modelowaniu wzajemnych relacji przestrzennych tych fragmentów, które pozostają widoczne. Wykorzystanie sekwencyjnej analizy map cech pozwala bowiem na utrzymanie spójności reprezentacji tożsamości pomimo braku ciągłości wizualnej w obrazie wejściowym.
\subsection{Dyskryminacyjność cech i obsługa próbek trudnych w podejściu opartym na marginesach kątowych}

\paragraph{}Alternatywnym podejściem do problemu okluzji jest potraktowanie go nie jako błędu wymagającego naprawy (rekonstrukcji), lecz jako specyficznego przypadku "trudnej próbki" (ang. hard sample), która generuje dużą wariancję wewnątrzklasową. Wiodące badania w tym zakresie koncentrują się na modyfikacji funkcji strat w taki sposób, aby wymusić większą separowalność klas w przestrzeni cech, co pozwala modelowi zachować skuteczność nawet przy niekompletnych danych wejściowych.

Deng i in. [3] w swojej pracy nad metodą ArcFace (Additive Angular Margin Loss) wykazują, że standardowa funkcja straty Softmax nie jest wystarczająca do radzenia sobie z dużymi zmianami wyglądu twarzy (w tym okluzją), ponieważ nie optymalizuje bezpośrednio podobieństwa wewnątrzklasowego. Wprowadzenie addytywnego marginesu kątowego pozwala na rzutowanie cech na hipersferę, co znacząco zwiększa dyskryminacyjność modelu. Co jednak kluczowe dla problemu okluzji, autorzy zauważają, że zaszumione lub przysłonięte twarze mogą zaburzać proces uczenia, jeśli wymusza się ich dopasowanie do pojedynczego centrum klasy.

\paragraph{}Jako rozwiązanie tego problemu Deng i in. proponują mechanizm Sub-center ArcFace. W tym podejściu każda tożsamość może posiadać $K$ pod-centrów (ang. sub-centers), a próbka ucząca musi znajdować się blisko tylko jednego z nich. Badania wykazały, że mechanizm ten prowadzi do automatycznej separacji danych: "czyste" i wyraźne twarze grupują się wokół dominującego pod-centrum, podczas gdy twarze przysłonięte, ustawione pod dużym kątem lub zaszumione, trafiają do pod-centrów niedominujących [3].

Z perspektywy implementacji systemów odpornych na okluzję, wnioski płynące z analizy ArcFace sugerują, że:

1. Zwiększanie marginesu geometrycznego w funkcji straty jest skuteczną metodą poprawy odporności na brakujące cechy.

2. Okluzja może być modelowana jako osobna modalność w ramach tej samej tożsamości (poprzez sub-centers), co zapobiega degradacji głównego wzorca twarzy podczas treningu sieci.
\subsection{Redukcja zakłóceń i adaptacyjna selekcja cech za pomocą modułów atencji (CBAM)}

\paragraph{}W kontekście rozpoznawania twarzy z okluzją kluczowym problemem jest nie tylko ekstrakcja cech, ale także umiejętność odróżnienia cech użytecznych (np. oczy, nos) od zakłóceń (np. szalik, okulary, dłoń zasłaniająca twarz). Rozwiązaniem tego problemu jest zastosowanie mechanizmów atencji, które pozwalają sieci neuronowej dynamicznie ważyć ważność poszczególnych fragmentów obrazu.

Woo i in. [4] zaproponowali moduł CBAM (Convolutional Block Attention Module), który jest lekkim i uniwersalnym rozszerzeniem dla sieci splotowych. Istotą tego rozwiązania jest sekwencyjne wnioskowanie map atencji w dwóch wymiarach:

1. Uwaga kanałowa (Channel Attention): Określa "co" jest ważne na obrazie. Wykorzystuje ona zarówno average-pooling (dla ogólnych statystyk), jak i max-pooling (dla cech wyróżniających), aby zdecydować, które kanały mapy cech zawierają istotne informacje o obiekcie, a które niosą szum.

2. Uwaga przestrzenna (Spatial Attention): Określa "gdzie" znajdują się istotne informacje. Jest ona komplementarna do atencji kanałowej i pozwala sieci skupić się na lokalizacji kluczowych elementów twarzy, ignorując tło i elementy przysłaniające.

\paragraph{}Badania przeprowadzone przez autorów, m.in. z wykorzystaniem techniki wizualizacji Grad-CAM, wykazały, że sieci wzbogacone o moduł CBAM znacznie precyzyjniej obejmują obszar docelowego obiektu. Woo i in. zauważają, że moduł ten prowadzi do redukcji szumu pochodzącego z nieistotnych elementów otoczenia (ang. clutters) [4]. W przypadku okluzji jest to właściwość fundamentalna – element przysłaniający twarz może zostać potraktowany przez sieć jako "nieistotny szum" (clutter) i zostać wytłumiony na etapie atencji przestrzennej, podczas gdy widoczne fragmenty twarzy zostaną wzmocnione.

Zastosowanie CBAM jest szczególnie obiecujące w implementacji modeli odpornych na okluzję ze względu na jego modułowość. Może on być zintegrowany z dowolną architekturą CNN (np. ResNet) przy znikomym narzucie obliczeniowym, co czyni go dogodnym narzędziem do analizy wpływu selekcji cech na różnych poziomach reprezentacji.
\subsection{Przegląd technik głębokiego uczenia i wyzwania związane ze zbiorami danych w analizie twarzy przysłoniętych}

\paragraph{}Kompleksowe ujęcie problemu analizy twarzy z okluzją wymaga uwzględnienia nie tylko architektury sieci, ale także specyfiki danych treningowych i systematyki stosowanych metod. W najnowszym artykule przeglądowym z 2025 roku, Thaher i in. [5] wskazują, że okluzja stanowi jedno z najważniejszych otwartych wyzwań dla systemów biometrycznych i nadzoru wizyjnego, drastycznie wpływając na skuteczność algorytmów detekcji i rozpoznawania.

Autorzy dokonują systematyzacji podejść badawczych, wyróżniając metody oparte na głębokim uczeniu jako dominujący standard. Wskazują przy tym na dwa główne nurty radzenia sobie z problemem:

1. Metody oparte na danych (Data-centric): Polegające na augmentacji zbiorów treningowych poprzez syntetyczne nakładanie masek lub innych obiektów przysłaniających na twarze, co pozwala modelom "nauczyć się" okluzji.

2. Metody oparte na modelu (Model-centric): Obejmujące projektowanie specyficznych architektur, takich jak sieci wykorzystujące mechanizmy uwagi (np. wspomniany wcześniej CBAM) czy piramidy cech (FPN), które pozwalają na ekstrakcję informacji z widocznych fragmentów twarzy mimo zakłóceń.

\paragraph{}Istotnym aspektem poruszonym w przeglądzie jest rola specjalistycznych zbiorów danych. Thaher i in. podkreślają, że standardowe bazy twarzy są niewystarczające do trenowania odpornych systemów. Wymieniają kluczowe zbiory dedykowane problemowi okluzji, takie jak MAFA (Masked Faces), które zawierają twarze z naturalnymi przesłonięciami, a nie tylko syntetycznymi. Autorzy zwracają uwagę, że sukces w implementacji modeli odpornych na okluzję zależy w dużej mierze od jakości i różnorodności tych danych, a także od zdolności algorytmów do radzenia sobie z ekstremalnymi warunkami, takimi jak połączenie okluzji z trudnym oświetleniem czy nietypową pozą [5].

Wnioski z tego przeglądu uzasadniają konieczność testowania implementowanych rozwiązań nie tylko na ogólnych zbiorach walidacyjnych, ale także na dedykowanych podzbiorach zawierających "twardą" okluzję, aby rzetelnie ocenić ich przydatność w rzeczywistych zastosowaniach.

Pomimo intensywnych badań nad rozpoznawaniem twarzy w warunkach okluzji, literatura wciąż w ograniczonym stopniu analizuje wpływ lekkich, modułowych mechanizmów atencji na odporność nowoczesnych architektur opartych na funkcjach strat z marginesem kątowym. W szczególności brakuje jednoznacznych analiz dotyczących umiejscowienia takich mechanizmów w strukturze sieci oraz ich wpływu na skuteczność identyfikacji w warunkach syntetycznej i rzeczywistej okluzji.

\section{Cel i zakres projektu}
\subsection{Cel projektu}

\paragraph{}Celem niniejszego projektu jest analiza wpływu częściowej okluzji twarzy, ze szczególnym uwzględnieniem zasłonięcia obszaru oczu, na skuteczność systemów rozpoznawania tożsamości w zadaniu identyfikacji osób. Praca koncentruje się na ocenie odporności współczesnych modeli głębokiego uczenia na zakłócenia wizualne typowe dla rzeczywistych scenariuszy, takich jak monitoring wizyjny czy zastosowania mobilne.

W ramach projektu zaprojektowano i zaimplementowano modyfikacje architektury bazowej IResNet50 (ArcFace), polegające na integracji mechanizmu uwagi typu Convolutional Block Attention Module (CBAM) w blokach rezydualnych oraz na zastosowaniu pomocniczej gałęzi Transformer jako dodatkowej funkcji straty. Celem tych działań było zwiększenie zdolności modeli do selektywnej ekstrakcji cech dyskryminacyjnych w warunkach częściowej utraty informacji wizualnej.

Skuteczność zaproponowanych rozwiązań oceniono poprzez porównanie z wybranymi modelami typu state-of-the-art, w tym zarówno ciężkimi architekturami referencyjnymi, jak i ich lekkimi wariantami przeznaczonymi do zastosowań o ograniczonych zasobach obliczeniowych. Przeprowadzone eksperymenty umożliwiają sformułowanie wniosków dotyczących wpływu okluzji na proces identyfikacji oraz efektywności architektonicznych modyfikacji w kontekście zwiększania odporności modeli na tego typu zakłócenia.
\subsection{Zakres pracy}

\paragraph{}Zakres pracy obejmuje analizę problemu rozpoznawania twarzy na podstawie obrazów 2D w warunkach częściowej okluzji, ze szczególnym uwzględnieniem syntetycznego zasłonięcia obszaru oczu. Badania przeprowadzono z wykorzystaniem wybranego zbioru danych zawierającego obrazy twarzy poddane procedurze wyrównania (alignment) oraz kontrolowanej augmentacji w postaci okluzji.

\paragraph{}W ramach projektu wykonano:

* porównanie wybranych modeli typu state-of-the-art (ArcFace Large, ArcFace Small, VGGFace) z autorskimi implementacjami opartymi na architekturze IResNet50,

* analizę wpływu rozproszonej uwagi CBAM (integracja w blokach rezydualnych) na skuteczność identyfikacji,

* analizę wpływu pomocniczej gałęzi Transformer na odporność na okluzję,

* ocenę skuteczności strategii transfer learningu i fine-tuningu w kontekście odporności na okluzję,

* ewaluację modeli w zadaniu identyfikacji tożsamości z wykorzystaniem metryk Rank-1 oraz Rank-3.

Zakres pracy ograniczono do porównania modelu bazowego oraz jego modyfikowanych wariantów architektonicznych. Nie przeprowadzono analizy alternatywnych rodzin modeli (np. w pełni transformatorowych), metod wymagających dodatkowych modalności biometrycznych ani zaawansowanych technik rekonstrukcji okluzji.

\paragraph{}Projekt nie obejmuje analizy danych wideo, danych trójwymiarowych (3D), zagadnień związanych z detekcją twarzy, ani optymalizacji modeli pod kątem wdrożeń produkcyjnych w systemach czasu rzeczywistego. Detekcja twarzy została wykorzystana wyłącznie jako etap wstępnego wyrównania danych wejściowych i nie stanowi przedmiotu analizy ani porównań w niniejszej pracy. Ponadto badania ograniczono do eksperymentów z wykorzystaniem syntetycznie generowanej okluzji, bez walidacji na rzeczywistych zbiorach danych zawierających naturalne przesłonięcia twarzy.

\section{Metodologia}
Niniejsza sekcja opisuje metodologię zastosowaną w opracowaniu i ocenie modeli rozpoznawania twarzy z naciskiem na niezawodność wobec okluzji. Szczegółowo omawiamy przygotowanie danych, architektury modeli, procedury treningowe i protokoły ewaluacji.

\subsection{Przygotowanie danych i przetwarzanie wstępne}
\subsubsection{Przegląd datasetu}

\paragraph{}Badania przeprowadzono na zbiorze danych CASIA-WebFace, zawierającym ponad 494 tys. kolorowych fotografii twarzy, należących do 10,575 unikatowych tożsamości. Dane podzielono na dwa rozłączne podzbiory: treningowy (90\%) i testowy (10\%). Dla monitorowania zbieżności modelu wykorzystywano stały zestaw 500 par (walidacja) z części testowej, zgodnie z pipeline’em treningowym.

\subsubsection{Wyrównanie na podstawie punktów orientacyjnych}

\paragraph{}Wszystkie obrazy są wstępnie przetwarzane za pomocą modelu RetinaFace od InsightFace w celu wykrycia 5 punktów orientacyjnych twarzy: lewe oko, prawe oko, czubek nosa i kąciki ust. Te punkty orientacyjne umożliwiają konsystentną normalizację twarzy przy użyciu transformacji podobieństwa do standardowej rozdzielczości 112×112. Większość obrazów ma towarzyszący plik JSON zawierający:
\begin{itemize}
    \item Współrzędne ramki ograniczającej: \texttt{bbox} = $[x_1, y_1, x_2, y_2]$
    \item 5 punktów orientacyjnych: \texttt{landmarks} = $\{$lewe\_oko, prawe\_oko, nos, lewy\_ust, prawy\_ust$\}$
\end{itemize}

W przypadku braku JSON, zastosowano rezerwowe skalowanie obrazu do 112×112. Brak pliku JSON wynikał z niepowodzenia detekcji punktów orientacyjnych dla niskiej jakości lub nietypowych obrazów (ok. 2,5\% zdjęć).

Wyrównanie twarzy jest wykonywane przy użyciu referencyjnej konfiguracji punktów orientacyjnych ArcFace:
\begin{equation}
M = \text{Similarity Transform}(\text{wykryte\_punkty}, \text{punkty\_referencyjne})
\end{equation}

gdzie punkty referencyjne dla obrazów 112×112 są ustalone na:
\begin{equation}
\text{arcface\_dest} = \begin{bmatrix} 38.29 & 51.70 \\ 73.53 & 51.50 \\ 56.03 & 71.74 \\ 41.55 & 92.37 \\ 70.73 & 92.20 \end{bmatrix}
\end{equation}

Punkty te stanowią standardową konfigurację stosowaną w ArcFace i zostały zaczerpnięte z dokumentacji kodu oryginalnego modelu, gdzie określono je jako referencyjne pozycje cech twarzy w wyrównanych obrazach.

Ta procedura zapewniła konsystentne wyrównanie zgodne z oryginalnym modelem Buffalo\_L (ArcFace Large), co jest kluczowe dla transferu wag.

\subsubsection{Strategia augmentacji danych}
\label{subsec:augmentation}

Wszystkie modele są trenowane z syntetyczną okluzją w celu symulacji zasłonięć w regionie oczu. Implementacja stosuje poziomy pasek okluzji o stałej wysokości i losowym kolorze.

\begin{itemize}
    \item \textbf{Wysokość}: 20 pikseli
    \item \textbf{Szerokość}: pełna szerokość obrazu (0 do 112)
    \item \textbf{Pozycja pionowa}: $y = 52 \pm 5$ pikseli (region oczu)
    \item \textbf{Kolor}: losowe RGB $(r, g, b)$, $r,g,b \in [0, 255]$
\end{itemize}

Prawdopodobieństwo wystąpienia okluzji różni się między wariantami zgodnie z kodem treningowym:
\begin{itemize}
    \item \textbf{Aligned\_Pretrained}: $p=0.7$
    \item \textbf{Aligned\_Pretrained\_Transformers}: $p=0.7$
    \item \textbf{CBAM\_Block}: $p=0.5$
\end{itemize}

W fazie testowej kolor okluzji jest stały (czarny), co zapewnia powtarzalność i porównywalność wyników pomiędzy modelami.

\subsection{Ewaluacja modeli typu SOTA}
W celu zapewnienia rzetelnego punktu odniesienia oraz właściwej kontekstualizacji uzyskanych wyników przeprowadzono ewaluację wybranych modeli typu state-of-the-art (SOTA), w tym modeli dostępnych w bibliotece InsightFace oraz modelu referencyjnego.

Wybrane modele obejmowały zarówno rozwiązania o wysokiej dokładności, jak i lekkie architektury zoptymalizowane pod kątem wydajności obliczeniowej, a także model referencyjny starszej generacji.

Testowane modele:
\begin{itemize}
\item ArcFace Large (buffalo\_l) – model oparty na architekturze IResNet-50, dostępny w bibliotece InsightFace,
\item ArcFace Small (buffalo\_s) – lekka architektura oparta na MobileFaceNet, przeznaczona do zastosowań mobilnych, dostępna w bibliotece InsightFace,
\item VGGFace2 (ResNet-50) – model referencyjny starszej generacji (2018 r.), zaimplementowany z wykorzystaniem biblioteki \textit{keras-vggface}, niebędący częścią frameworka InsightFace.
\end{itemize}

Wszystkie modele poddano temu samemu protokołowi identyfikacji (galeria–sonda) oraz symulacji okluzji w postaci poziomego zasłonięcia regionu oczu o wysokości 20 pikseli. Należy jednak podkreślić, że część modeli SOTA wykorzystuje własny pipeline ekstrakcji cech: dla ArcFace (InsightFace) embeddingi są wyznaczane na podstawie pełnego obrazu z detekcją twarzy, natomiast dla VGGFace (ResNet50) stosowane jest przycięcie względem \texttt{bbox} i skalowanie do 224×224. Zostało to utrzymane zgodnie z implementacjami referencyjnymi w repozytorium. Szczegółowy opis protokołu ewaluacji przedstawiono w dalszej części dokumentu.

\subsection{Model bazowy}
\subsubsection{Wybór IResNet50}
Wybraliśmy IResNet50 (Ulepszona sieć ResNet50) jako architekturę sieci głównej ze względu na jej sprawdzoną wydajność w modelu ArcFace-Large (buffalo\_l) z projektu InsightFace. Implementacja w repozytorium zawiera warstwy BN i PReLU, a warstwa \texttt{features} (BN1d) ma zamrożone wagi (\texttt{requires\_grad=False}), co ogranicza adaptację końcowego wektora cech w trakcie fine-tuningu. IResNet50 składa się z:
\begin{itemize}
    \item Wstępna konwolucja: $3 \times 3$ conv, 64 kanały, stride 1
    \item Warstwy 1-4: Bloki resztkowe o różnych głębokościach
    \item Normalizacja wsadowa i aktywacje PReLU na całej sieci
    \item Wymiar cechy: osadzania 512-wymiarowe
\end{itemize}

\subsubsection{Hybrydowa strategia transferu wag}
Bezpośredni transfer z modelu buffalo\_l opartego na ONNX do PyTorch okazał się trudny ze względu na niezgodności architektoniczne. Zastosowano hybrydowe dopasowanie wag (dopasowanie nazw + dopasowanie kształtu), z reinicjalizacją warstw bez dopasowania. Proces ten jest realizowany przez \texttt{load.py} i zapewnia stabilny punkt startowy dla fine-tuningu.

\paragraph{Pomyślne transfery:} $\approx 90\%$ wag bloków ResNet zostało pomyślnie zmapowanych, w tym:
\begin{itemize}
    \item Warstwy Conv2d w ścieżkach resztkowych
    \item Większość parametrów BatchNorm (wagi, obciążenia, statystyki uruchamiające)
    \item Wagi PReLU (gdzie konwencje nazewnictwa się zgadzały)
\end{itemize}

\paragraph{Warstwy reinicjalizowane:} 26 warstw wymagało inicjalizacji Kaiming Normal (He) ze względu na niezgodność architektoniczną:
\begin{itemize}
    \item layer1.2.bn3: statystyki BatchNorm3D (running\_mean, running\_var)
    \item layer2.3.bn2, layer2.3.bn3: parametry BatchNorm
    \item layer2.3.prelu: waga PReLU (parametr $\alpha$)
    \item layer3.12-13.bn2-3: podobne instancje BatchNorm i PReLU
\end{itemize}

Formuła inicjalizacji:
\begin{equation}
w \sim \mathcal{N}(0, \sqrt{\frac{2}{n_{in}(1 + a^2)}})
\end{equation}
gdzie $a$ jest ujemnym nachyleniem (0 dla ReLU, 0.25 dla PReLU w naszym przypadku - później trenowane).

\subsection{Funkcje straty}
\label{subsec:loss}

\subsubsection{Strata ArcMargin}
Podstawową funkcją straty w modelu bazowym oraz w wariancie CBAM\_Block jest strata ArcMargin, marginalna strata softmax, która wymusza separowalność kątową w przestrzeni osadzania. W implementacji przyjęto parametry $m=0.5$ i $s=30.0$ (zbieżne z ArcFace).

W treningu optymalizowana jest standardowa \textit{CrossEntropy} na wyjściu ArcMargin:
\begin{equation}
\mathcal{L}_{\text{ArcFace}} = \text{CE}(\text{ArcMargin}(\mathbf{e}), y)
\end{equation}

\subsubsection{Strata złożona dla wariantu Transformer}
W wariancie Transformer zastosowano funkcję straty łączoną, gdzie gałąź podstawowa (ArcMargin) jest wzbogacana o pomocniczą stratę klasyfikacji z transformera (CrossEntropy):
\begin{equation}
\mathcal{L}_{\text{total}} = (1-\alpha)\,\mathcal{L}_{\text{ArcFace}} + \alpha\,\mathcal{L}_{\text{Trans}}
\end{equation}
gdzie $\alpha = 0.4$ zgodnie z parametrami w kodzie treningu.

\subsection{Modyfikacja Architektury: Strategie Odporności na Okluzję}
Zaproponowane modyfikacje architektury ograniczono do wariantów faktycznie obecnych w implementacji. Wszystkie modele bazują na IResNet50 i różnią się wyłącznie dodatkowymi mechanizmami odporności na okluzję.

\paragraph{Wariant A: Model bazowy}
Wariant odniesienia wykorzystuje standardowy backbone IResNet50 oraz głowicę ArcMargin do klasyfikacji tożsamości. Nie stosuje mechanizmów uwagi ani dodatkowych głowic pomocniczych. Jego rolą jest stanowić stabilny punkt odniesienia do oceny wpływu modyfikacji architektury.

\paragraph{Wariant B: CBAM zintegrowany w blokach rezydualnych}
W tym wariancie moduły CBAM umieszczono wewnątrz każdego bloku rezydualnego (CBAMBasicBlock), bezpośrednio przed połączeniem rezydualnym. Zastosowano uwagę kanałową i przestrzenną, co pozwala na selektywne wzmacnianie istotnych map cech na wielu poziomach abstrakcji. Celem wariantu jest uzyskanie odporności na okluzję poprzez rozproszoną filtrację zakłóceń.

\paragraph{Wariant C: Gałąź pomocnicza Transformer}
Wariant ten modyfikuje backbone tak, aby zwracał mapę cech o rozmiarze $7\times7$, a następnie dodaje pomocniczą gałąź Transformer (Branch-2) uczoną równolegle z głowicą ArcMargin (Branch-1). W trakcie treningu wykorzystywana jest strata łączona (ArcMargin + CrossEntropy z gałęzi Transformer), natomiast podczas ewaluacji wykorzystywana jest wyłącznie gałąź standardowa.

\subsubsection{Strategia uczenia i stabilizacji}
We wszystkich wariantach zastosowano transfer learning z wagami wstępnie wytrenowanymi, a w wariancie CBAM\_Block dodatkowo etapowe zamrażanie/odmrażanie wag w celu stabilizacji uczenia. Augmentacja danych obejmowała syntetyczną okluzję regionu oczu, spójną z protokołem ewaluacyjnym.


\subsection{Procedura Treningu}
\subsubsection{Model bazowy (Aligned\_Pretrained)}
Model bazowy trenowany jest jako klasyfikator z głowicą ArcMargin.

\begin{table}[h]
\centering
\small
\begin{tabular}{ll}
\toprule
\textbf{Hiperparametr} & \textbf{Wartość} \\
\midrule
Rozmiar partii & 32 \\
Epoki & 25 \\
Optymalizator & SGD (momentum=0.9) \\
LR & $1 \times 10^{-2}$ \\
Zanik wag & brak (nie ustawiono) \\
Wczesne zatrzymanie & Tak (cierpliwość=5) \\
Walidacja & 500 par z testu (średnia cos) \\
Augmentacja okluzji & $p=0.7$, wysokość 20 px \\
\bottomrule
\end{tabular}
\caption{Hiperparametry modelu bazowego (Aligned\_Pretrained).}
\label{tab:baseline_config}
\end{table}

\subsubsection{Gałąź pomocnicza Transformer}
Model z transformatorem wykorzystuje złożoną funkcję straty, łącząc ArcMargin i klasyfikację gałęzi pomocniczej.

\begin{table}[h]
\centering
\small
\begin{tabular}{ll}
\toprule
\textbf{Hiperparametr} & \textbf{Wartość} \\
\midrule
Rozmiar partii & 32 \\
Epoki & 25 \\
Optymalizator & SGD (momentum=0.9, weight\_decay=$5\times 10^{-4}$) \\
LR & $1 \times 10^{-2}$ \\
ALPHA (balans strat) & 0.4 \\
Transformer layers/heads & 6 / 8 \\
Wczesne zatrzymanie & Tak (cierpliwość=5) \\
Walidacja & 500 par z testu (średnia cos) \\
Augmentacja okluzji & $p=0.7$, wysokość 20 px \\
\bottomrule
\end{tabular}
\caption{Hiperparametry modelu z gałęzią Transformer.}
\label{tab:transformer_config}
\end{table}

\subsubsection{CBAM\_Block (uwaga w każdym bloku)}
Trening realizowany jest dwuetapowo z zamrożeniem backbone'u w pierwszej epoce.

\begin{table}[h]
\centering
\small
\begin{tabular}{ll}
\toprule
\textbf{Hiperparametr} & \textbf{Wartość} \\
\midrule
Rozmiar partii & 64 \\
Epoki & 30 \\
Optymalizator & SGD (momentum=0.9, weight\_decay=$5\times 10^{-4}$) \\
LR start & 0.1 \\
Harmonogram LR & StepLR (step\_size=10, gamma=0.1) \\
Zamrażanie & Epoka 1: tylko CBAM + head \\
Walidacja & 500 par z testu (średnia cos) \\
Augmentacja okluzji & $p=0.5$, wysokość 20 px \\
\bottomrule
\end{tabular}
\caption{Hiperparametry modelu CBAM\_Block.}
\label{tab:cbam_block_config}
\end{table}

\subsection{Protokół ewaluacji}
\subsubsection{Konstrukcja zestawu testowego}
Protokół ewaluacji stosuje zadanie identyfikacji galeria-vs-sonda:

\paragraph{Zestaw galerii:} Skonstruowany z pierwszej połowy sesji (folderów) dla każdej osoby w podziale testowym. Dla każdej osoby osadzania z wielu czystych (nie okluzjowanych) obrazów są uśredniane:

\begin{equation}
\mathbf{e}_{\text{gallery}} = \frac{1}{n} \sum_{i=1}^{n} \text{L2-Normalize}(\mathbf{e}_i)
\end{equation}

Zapewnia to niezawodny prototyp osoby do dopasowania podobieństwa.

\paragraph{Zestaw sondy:} Skonstruowany z drugiej połowy sesji. Te obrazy są poddane syntetycznej okluzji podczas ewaluacji (szczegóły w sekcji \ref{subsubsec:eval_occlusion}).

\subsubsection{Ekstrakcja embeddingu i indeksowanie}
\begin{enumerate}
    \item Ekstrahuj embedding dla wszystkich obrazów galerii i sondy
    \item Normalizuj L2: $\hat{\mathbf{e}} = \frac{\mathbf{e}}{\|\mathbf{e}\|_2}$
    \item Indeksuj embedding galerii przy użyciu FAISS IndexFlatIP (Inner Product), gdzie podobieństwo cosinusowe na wektorach znormalizowanych L2 równa się iloczynowi wewnętrznemu
\end{enumerate}

\paragraph{Rozdział walidacji i ewaluacji:} w trakcie treningu stosowana jest pomocnicza walidacja na 500 parach (średnia cosinusowa czyste vs okluzjowane), próbkowanych z części testowej. Wyniki Rank-1/Rank-3 pochodzą z niezależnego pipeline’u FAISS dla zadania identyfikacji galeria-vs-sonda.

\subsubsection{Symulacja okluzji w ewaluacji}
\label{subsubsec:eval_occlusion}

Podczas ewaluacji każdy obraz sondy jest syntetycznie okluzjowany przy użyciu deterministycznego protokołu, aby zapewnić powtarzalność:

\begin{itemize}
    \item \textbf{Wysokość okluzji}: 20 pikseli
    \item \textbf{Szerokość okluzji}: jeśli dostępny \texttt{bbox} z adnotacji, pasek ograniczany jest do zakresu $[x_1, x_2]$; w przeciwnym razie pełna szerokość obrazu (0 do 112)
    \item \textbf{Kolor}: Czarny (0,0,0) [deterministyczny, w przeciwieństwie do losowej augmentacji treningowej]
    \item \textbf{Pozycja pionowa}: Region oczu, obliczony z punktów orientacyjnych:
    \begin{equation}
    c_y = \frac{y_{\text{lewe\_oko}} + y_{\text{prawe\_oko}}}{2}
    \end{equation}
    \item \textbf{Fallback}: Jeśli punkty orientacyjne niedostępne, $c_y = \frac{h}{2} - 10 = 46$ (dla 112×112)
    \item \textbf{Zakres Y}: $[c_y - 10, c_y + 10]$
\end{itemize}

Deterministyczna czarna okluzja (w przeciwieństwie do losowych kolorów w treningu) zapewnia rygorystyczną, powtarzalną ewaluację podczas testowania uogólnienia na wariacje kolorów.

\paragraph{Wyrównanie w ewaluacji:} obrazy testowe w zbiorze \texttt{webface\_112x112} są już wstępnie wyrównane.

\subsubsection{Wyszukiwanie podobieństwa i metryki}
Dla każdego obrazu sondy $p$:
\begin{enumerate}
    \item Zapytaj indeks FAISS z $k=3$ najbliższymi sąsiadami
    \item Pobierz top-3 tożsamości galerii ze wynikami podobieństwa: $(s_1, s_2, s_3)$
\end{enumerate}

\paragraph{Dokładność Rank-k:} 
\begin{equation}
\text{Rank-}k = \frac{\text{\# zapytań z prawidłową tożsamością w top-}k}{N_{\text{ogółem}} }
\end{equation}

W szczególności:
\begin{itemize}
    \item \textbf{Dokładność Rank-1}: $P(\text{tożsamość}[0] = \text{rzeczywistość})$
    \item \textbf{Dokładność Rank-3}: $P(\text{rzeczywistość} \in \{\text{tożsamość}[0,1,2]\})$
\end{itemize}

\subsubsection{Wyjście i logowanie}
Szczegółowe wyniki są logowane w formacie CSV z kolumnami:
\begin{lstlisting}
query_id, top1_id, sim1, top2_id, sim2, top3_id, sim3, found_in_top3
\end{lstlisting}

\subsection{Środowisko obliczeniowe i implementacja}
\paragraph{}Wszystkie eksperymenty zostały przeprowadzone w środowisku Python z wykorzystaniem PyTorch do trenowania modeli, InsightFace do detekcji punktów orientacyjnych oraz FAISS do szybkiego wyszukiwania podobieństwa w przestrzeni embedding'ów. Dla modeli autorskich zastosowano wspólny pipeline danych (wyrównanie, augmentacja, normalizacja, ekstrakcja embedding'ów), co minimalizuje ryzyko różnic wynikających z niejednorodnych procedur wejściowych. Modele SOTA korzystały z własnych, referencyjnych pipeline’ów zgodnie z implementacjami w repozytorium.

Modelowanie architektur odbywało się w tym samym szkielecie IResNet50, a modyfikacje ograniczano do ściśle kontrolowanych miejsc (gałąź pomocnicza Transformer lub moduły CBAM). Taka polityka zmian pozwala przypisać obserwowane różnice wydajności wyłącznie wprowadzonym mechanizmom odporności na okluzję.

\subsection{Dobór hiperparametrów i stabilność uczenia}
\paragraph{}Wybór hiperparametrów oparto na konserwatywnym dostrajaniu: najpierw ustalono stabilny punkt odniesienia (baseline), następnie wprowadzono modyfikacje architektury (CBAM\_Block i Transformer) przy zachowaniu porównywalnego pipeline’u danych. Dla wariantu CBAM\_Block stosowano harmonogram StepLR oraz etapowe zamrażanie/odmrażanie warstw, a dla pozostałych wariantów wczesne zatrzymanie na podstawie walidacji (500 par) w celu ograniczenia niestabilności numerycznej.

Parametry uczenia (rozmiar partii, epoki, współczynniki uczenia) dobrano empirycznie na podstawie obserwowanej zbieżności funkcji straty i stabilności średniej cosinusowej na zbiorze walidacyjnym (500 par z części testowej). Jednolity protokół augmentacji (okluzja w rejonie oczu) utrzymano we wszystkich wariantach, co umożliwia bezpośrednie porównanie modeli.

\subsection{Reprodukowalność i logowanie}
\paragraph{}W kodzie nie ustawiono stałych ziaren losowości, dlatego przebiegi mogą wykazywać niewielką wariancję między uruchomieniami. Dla każdego wariantu zachowywano najlepszy checkpoint według średniej cosinusowej na walidacji (500 par z części testowej), a końcowa ewaluacja wykonywana była wyłącznie na ustalonym zbiorze testowym z deterministyczną okluzją. Wyniki zbiorcze zostały zapisane w postaci tabel zbiorczych, co ułatwia replikację i weryfikację rezultatów.

\section{Wyniki}
\paragraph{}Niniejsza sekcja prezentuje wyniki ewaluacji modeli rozpoznawania twarzy na testowym zbiorze CASIA-WebFace, ze szczególnym uwzględnieniem ich odporności na okluzję. Wartości Rank-1/Rank-3 pochodzą z zapisów \texttt{score.txt} wygenerowanych przez pipeline FAISS w scenariuszu okluzji 20px. Wyniki zawierają porównanie dokładności Rank-k, analizę wpływu syntetycznej okluzji oraz porównanie z modelami SOTA.

\subsection{Porównanie z modelami SOTA}

\paragraph{}W celu ustalenia benchmarku, poddaliśmy testowaniu trzy modele stanowiące stan techniki w rozpoznawaniu twarzy: ArcFace Large (buffalo\_l), ArcFace Small (buffalo\_s) oraz VGGFace (model legacy). Wszystkie modele testowano w identycznym scenariuszu okluzji (syntetyczna okluzja 20 pikseli na poziomie oczu). Wyniki zawarto w Tabeli~\ref{tab:sota_baseline}.

\begin{table}[H]
\centering
\small
\begin{tabular}{lcc}
\toprule
\textbf{Model} & \textbf{Rank-1 (\%)} & \textbf{Rank-3 (\%)} \\
\midrule
ArcFace Large (buffalo\_l) & 94.81 & 95.51 \\
ArcFace Small (buffalo\_s) & 86.65 & 90.68 \\
VGGFace (Legacy) & 39.64 & 57.01 \\
\bottomrule
\end{tabular}
\caption{Wyniki modeli SOTA na zbiorze testowym z syntetyczną okluzją 20px.}
\label{tab:sota_baseline}
\end{table}

\paragraph{}ArcFace Large wyznacza górny limit wydajności (94.81\%), przy czym model wykorzystuje masywny dataset treningowy (MS1MV2 z $\sim$5.8M obrazów). ArcFace Small (86.65\%) stanowi punkt odniesienia dla efektywnych modeli przemysłowych. VGGFace (39.64\%) pokazuje drastyczną degradację, podkreślając znaczenie nowoczesnych funkcji straty opartych na marginesach (ArcFace) w porównaniu do starszych podejść opartych na softmax.

\subsection{Dokładność modeli badawczych z okluzją}

\paragraph{}Modele własne obecne w repozytorium obejmują wariant bazowy, wariant z CBAM w blokach rezydualnych oraz wariant z gałęzią Transformer. Wyniki Rank-1 i Rank-3 przedstawiono w Tabeli~\ref{tab:results_custom}.

\begin{table}[H]
\centering
\small
\begin{tabular}{lccc}
\toprule
\textbf{Model} & \textbf{Rank-1 (\%)} & \textbf{Rank-3 (\%)} & \textbf{Delta vs Baseline} \\
\midrule
Baseline & 87.11 & 90.56 & -- \\
CBAM\_Block & 90.16 & 91.88 & +3.05 pp \\
Transformers & 85.00 & 89.02 & -2.11 pp \\
\bottomrule
\end{tabular}
\caption{Dokładność modeli badawczych na zbiorze testowym z syntetyczną okluzją 20px.}
\label{tab:results_custom}
\end{table}

\paragraph{}Najlepsze wyniki wśród modeli badawczych osiągnął wariant CBAM\_Block (90.16\%), przekraczając baseline o 3.05 pp. Wariant Transformer uzyskał niższy wynik niż baseline.

\subsection{Analiza degradacji wydajności}

\paragraph{}Porównanie wyników modeli w scenariuszu okluzji, nieznaczenie, ale potwierdza, że rozproszona uwaga (CBAM\_Block) zwiększa odporność na zasłonięcie kluczowych regionów. Model ArcFace Large osiąga 94.81\% w warunkach okluzji 20px, podczas gdy najlepszy wariant badawczy CBAM\_Block uzyskuje 90.16\%.

\subsection{Ranking modeli i luka Rank-3}
\paragraph{}W celu analizy zachowania top-k zestawiono różnice Rank-3 vs Rank-1 (im mniejsza luka, tym większa pewność top-1). Najmniejszą lukę uzyskano dla ArcFace Large i CBAM\_Block, co wskazuje na stabilne dopasowania wśród najbliższych sąsiadów.

\begin{table}[H]
\centering
\small
\begin{tabular}{lccc}
\toprule
\textbf{Model} & \textbf{Rank-1 (\%)} & \textbf{Rank-3 (\%)} & \textbf{Luka Rank-3 - Rank-1} \\
\midrule
ArcFace Large (buffalo\_l) & 94.81 & 95.51 & 0.70 \\
CBAM\_Block & 90.16 & 91.88 & 1.72 \\
Baseline & 87.11 & 90.56 & 3.45 \\
ArcFace Small (buffalo\_s) & 86.65 & 90.68 & 4.03 \\
Transformers & 85.00 & 89.02 & 4.02 \\
VGGFace (Legacy) & 39.64 & 57.01 & 17.37 \\
\bottomrule
\end{tabular}
\caption{Luka pomiędzy Rank-3 i Rank-1 dla wszystkich modeli (mniejsza luka = większa pewność top-1).}
\label{tab:rank_gap}
\end{table}

\section{Dyskusja}
\paragraph{}W niniejszej sekcji interpretujemy uzyskane rezultaty w kontekście przyjętej metodologii, ze szczególnym uwzględnieniem czynników wpływających na trafność porównań między badanymi wariantami.

\subsection{Interpretacja wyników i analiza czynników konfundujących}
\paragraph{}W przeprowadzonych eksperymentach wariant CBAM\_Block osiągnął najwyższą dokładność Rank-1 (90.16\%), przewyższając model bazowy (87.11\%) o 3.05 pp. Należy jednak zaznaczyć, że wynik ten nie może być przypisany wyłącznie modyfikacji architektury, ze względu na różnice w procedurze treningowej.

Analiza konfiguracji modeli ujawnia istotne różnice w hiperparametrach, które mogły wpłynąć na wynik końcowy:
\begin{itemize}
	\item \textbf{Intensywność augmentacji:} Model bazowy był trenowany z prawdopodobieństwem okluzji $p=0.7$, podczas gdy wariant CBAM trenowano w warunkach łagodniejszych ($p=0.5$). Oznacza to, że model CBAM miał dostęp do większej liczby „czystych” próbek w trakcie uczenia, co naturalnie ułatwia zbieżność i może prowadzić do wyższych wyników na zbiorze testowym.
	\item \textbf{Dynamika uczenia (Learning Rate):} Wariant CBAM rozpoczął trening z 10-krotnie wyższym współczynnikiem uczenia ($LR=0.1$) w porównaniu do modelu bazowego ($LR=1\times10^{-2}$). Wyższa początkowa wartość LR, w połączeniu z harmonogramem StepLR, mogła pozwolić na lepszą eksplorację przestrzeni wag i ucieczkę z minimów lokalnych, niezależnie od obecności modułów uwagi.
\end{itemize}

W związku z powyższym, uzyskany wynik należy interpretować jako rezultat połączenia mechanizmu uwagi oraz bardziej agresywnej strategii optymalizacji przy mniejszym poziomie zakłóceń treningowych. Nie można jednoznacznie wyizolować wpływu samego modułu CBAM bez przeprowadzenia ponownych testów w warunkach ceteris paribus.

Wariant Transformer uzyskał wynik niższy od baseline (85.00\%), co przy zastosowaniu tej samej intensywności okluzji ($p=0.7$) sugeruje, że dodanie pomocniczej gałęzi i złożonej funkcji straty mogło utrudnić zbieżność modelu przy ograniczonej liczbie epok (25).

\subsection{Porównanie z modelami SOTA i kwestie wyrównania}
\paragraph{}Porównanie z modelami referencyjnymi (ArcFace Large: 94.81\%, ArcFace Small: 86.65\%) pełni rolę orientacyjną. Należy podkreślić różnice w potoku przetwarzania wstępnego (pre-processing). Modele autorskie korzystały z wyrównania opartego na detekcji RetinaFace do wymiaru $112\times112$, podczas gdy modele SOTA (szczególnie VGGFace) wykorzystują odmienne procedury kadrowania (crop) i skalowania. W biometrii twarzy sposób wyrównania ma krytyczny wpływ na skuteczność, dlatego bezpośrednie zestawienie wyników obarczone jest błędem systematycznym wynikającym z różnic w danych wejściowych.

Mimo to, fakt że wariant CBAM\_Block (90.16\%) przewyższył wynik ArcFace Small (86.65\%), wskazuje na potencjał zastosowanej konfiguracji (architektura + harmonogram treningu) w zastosowaniach o ograniczonych zasobach, nawet jeśli przewaga ta wynika częściowo z doboru hiperparametrów.

\subsection{Ograniczenia badania (Threats to Validity)}
\paragraph{}Krytyczna analiza przebiegu projektu wskazuje na następujące ograniczenia, które należy uwzględnić przy replikacji wyników:
\begin{itemize}
	\item \textbf{Niejednorodne warunki eksperymentalne:} Zmienność parametrów $p$ (prawdopodobieństwo okluzji) oraz $LR$ pomiędzy wariantami uniemożliwia precyzyjną ocenę przyrostu jakości wynikającego wyłącznie ze zmian w architekturze.
	\item \textbf{Jakość modelu bazowego (Baseline):} Procedura transferu wag wymagała losowej reinicjalizacji 26 warstw, w tym kluczowych statystyk Batch Normalization. Mogło to sztucznie obniżyć wydajność modelu bazowego, sprawiając, że stał się on zbyt łatwym punktem odniesienia dla modelu CBAM trenowanego od początku z wyższym LR.
	\item \textbf{Brak analizy statystycznej:} Ze względu na brak ustalonych ziaren losowości (seeds) oraz wykonanie pojedynczych przebiegów treningowych, obserwowane różnice mogą mieścić się w granicach błędu statystycznego.
\end{itemize}


\section{Wnioski}
\paragraph{}W ramach projektu przeprowadzono analizę wpływu modyfikacji architektury IResNet50 na skuteczność identyfikacji twarzy w warunkach okluzji. Na podstawie przeprowadzonych eksperymentów i krytycznej analizy wyników sformułowano następujące wnioski:

\begin{enumerate}
    \item \textbf{Wydajność wariantu z uwagą:} Wariant wykorzystujący moduły CBAM w blokach rezydualnych osiągnął najlepszy wynik w testach (Rank-1: 90.16\%), wykazując wysoką odporność na syntetyczną okluzję. Należy jednak zastrzec, że na sukces ten złożyła się nie tylko modyfikacja architektury, ale również zastosowanie odmiennej strategii uczenia (wyższe LR, łagodniejsza augmentacja) w porównaniu do modelu bazowego.
    \item \textbf{Wpływ gałęzi Transformer:} Wprowadzenie pomocniczej gałęzi Transformer przy zachowaniu silnej augmentacji ($p=0.7$) nie przyniosło poprawy wyników w zadanym reżimie treningowym (25 epok), co sugeruje, że bardziej złożone architektury wymagają dłuższego procesu uczenia lub większych zbiorów danych do poprawnej zbieżności.
    \item \textbf{Wyzwania metodologiczne:} Projekt uwypuklił kluczowe znaczenie rygorystycznej kontroli hiperparametrów (Learning Rate, Augmentation Probability) przy porównywaniu architektur sieci neuronowych. Zmienność tych parametrów pomiędzy eksperymentami została zidentyfikowana jako główny czynnik utrudniający jednoznaczną ocenę wpływu samych modułów uwagi.
    \item \textbf{Zalecenia do dalszych prac:} W celu weryfikacji postawionych hipotez, przyszłe prace powinny skupić się na wyrównaniu protokołu treningowego dla wszystkich modeli (strategia ceteris paribus), zastosowaniu walidacji krzyżowej oraz testach na realistycznych zbiorach okluzji (np. MAFA), aby wyeliminować wpływ czynników losowych i syntetycznego charakteru zakłóceń.
\end{enumerate}


\section{Bibliografia}
[1] - \textit{D. Zeng, R. Veldhuis i L. Spreeuwers, „A survey of face recognition techniques under occlusion”, arXiv preprint arXiv:2006.11366, 2020}

[2] - \textit{P. Prakash i A. K. Rai, „Age-Defying Face Recognition with Transformer-Enhanced Loss”, arXiv preprint arXiv:2412.02198, 2025}

[3] - \textit{J. Deng, J. Guo, J. Yang, N. Xue, I. Kotsia i S. Zafeiriou, „ArcFace: Additive Angular Margin Loss for Deep Face Recognition”, 2018}

[4] - \textit{S. Woo, J. Park, J.-Y. Lee i I. S. Kweon, „CBAM: Convolutional Block Attention Module”, arXiv preprint arXiv:1807.06521, 2018}

[5] - \textit{T. Thaher, M. Mafarja, M. Saffarini, A. H. H. M. Mohamed i A. A. El-Saleh, „A Comprehensive Review of Face Detection Techniques for Occluded Faces: Methods, Datasets, and Open Challenges”, Computer Modeling in Engineering \& Sciences, tom 144, nr 1, 2025.}

\begin{comment}
\end{comment}

\newpage
\appendix

\section{Szczegóły architektury CBAM}
\label{app:cbam}

\subsection{Moduł uwagi kanałów}
Moduł uwagi kanałów (CAM) oblicza wagi uwagi dla każdego kanału poprzez agregację wymiarów przestrzennych:

\subsubsection{Moduł uwagi kanałów}
Moduł uwagi kanałów (Channel Attention Module) generuje mapę wag $\mathbf{M}_c \in \mathbb{R}^{C \times 1 \times 1}$ poprzez eksploatację zależności międzykanałowych. Proces ten wykorzystuje współdzielony MLP operujący na zagregowanych przestrzennie cechach:

\begin{equation}
\mathbf{M}_c(\mathbf{F}) = \sigma\left( W_1(\text{ReLU}(W_0(\mathbf{F}^c_{\text{avg}}))) + W_1(\text{ReLU}(W_0(\mathbf{F}^c_{\text{max}}))) \right)
\end{equation}

gdzie:
\begin{itemize}
    \item $\mathbf{F}^c_{\text{avg}}, \mathbf{F}^c_{\text{max}}$ oznaczają wektory cech uzyskane odpowiednio przez Average Pooling i Max Pooling przestrzenny.
    \item $W_0 \in \mathbb{R}^{C/r \times C}$ oraz $W_1 \in \mathbb{R}^{C \times C/r}$ to wagi współdzielonej sieci MLP (bottleneck), a $r$ jest współczynnikiem redukcji.
    \item $\sigma$ oznacza funkcję aktywacji sigmoidalną.
\end{itemize}

Ostatecznie, mapa cech jest modyfikowana poprzez mnożenie elementarne (z broadcastowaniem wzdłuż wymiarów przestrzennych):

\begin{equation}
\mathbf{F}' = \mathbf{M}_c(\mathbf{F}) \otimes \mathbf{F}
\end{equation}

\subsection{Moduł uwagi przestrzennej}
Moduł uwagi przestrzennej (SAM) generuje mapy uwagi przestrzennej poprzez łączenie uśredniania kanałów i max poolingu:

\begin{equation}
\text{SAM}(\mathbf{X}) = \sigma(\text{Conv}_{7 \times 7}([\text{AvgPool}(\mathbf{X}); \text{MaxPool}(\mathbf{X})]))
\end{equation}

gdzie:
\begin{itemize}
    \item $[\cdot; \cdot]$ = łączenie kanałów (wejście: $h \times w \times 2$)
    \item $\text{Conv}_{7 \times 7}$ = pojedyncza konwolucja $7 \times 7$ z 1 kanałem wyjścia
    \item $\sigma$ = aktywacja sigmoidalna
\end{itemize}

Uwaga przestrzenna jest zastosowana jako:
\begin{equation}
\mathbf{X}'' = \text{SAM}(\mathbf{X}') \otimes \mathbf{X}'
\end{equation}

\subsection{Kompletny blok CBAM}
Kompletny blok CBAM łączy oba mechanizmy uwagi w sekwencji:

\begin{equation}
\begin{split}
    \mathbf{F}' &= \mathbf{M}_c(\mathbf{F}) \otimes \mathbf{F} \\
    \mathbf{F}'' &= \mathbf{M}_s(\mathbf{F}') \otimes \mathbf{F}'
\end{split}
\end{equation}

\begin{equation}
\mathbf{Y} = \mathbf{F}'' + \mathbf{X}
\end{equation}

W architekturach resztkowych CBAM jest stosowany przed operacją dodawania resztkowego:

\begin{equation}
\mathbf{Y} = \text{CBAM}(\text{Conv}(\mathbf{X})) + \mathbf{X}
\end{equation}

Dla IResNet50 zastosowano wariant \textbf{Block-wise (CBAM\_Block)}: moduł CBAM jest wstawiony wewnątrz każdego bloku resztkowego, co umożliwia hierarchiczną filtrację zakłóceń na wielu poziomach.

\section{Tabele uzupełniające}
\label{app:tables}

\subsection{Podsumowanie kompletnych konfiguracji modeli}

\begin{table}[H]
\centering
\scriptsize
\begin{tabular}{l|ccccc}
\toprule
\textbf{Model} & \textbf{Partia} & \textbf{Epoki} & \textbf{LR} & \textbf{Optimizer} & \textbf{Harmonogram} \\
\midrule
Aligned\_Pretrained & 32 & 25 & $1\times10^{-2}$ & SGD (mom=0.9) & -- \\
Aligned\_Pretrained\_Transformers & 32 & 25 & $1\times10^{-2}$ & SGD (mom=0.9, wd=$5\times10^{-4}$) & -- \\
CBAM\_Block & 64 & 30 & $0.1$ & SGD (mom=0.9, wd=$5\times10^{-4}$) & StepLR(10, 0.1) \\
\bottomrule
\end{tabular}
\caption{Podsumowanie konfiguracji treningowych modeli obecnych w repozytorium.}
\label{tab:all_configs}
\end{table}

\subsection{Mapowanie wariantów architektury na mechanizmy odporności}

\begin{table}[H]
\centering
\scriptsize
\begin{tabular}{lccc}
\toprule
\textbf{Wariant} & \textbf{Mechanizm} & \textbf{Miejsce integracji} & \textbf{Dodatkowa strata} \\
\midrule
Aligned\_Pretrained & Brak & -- & Nie \\
CBAM\_Block & Uwaga kanałowa + przestrzenna & W każdym bloku resztkowym & Nie \\
Aligned\_Pretrained\_Transformers & Gałąź Transformer & Po mapie cech $7\times7$ & Tak (CE, $\alpha=0.4$) \\
\bottomrule
\end{tabular}
\caption{Zestawienie wariantów architektury i mechanizmów odporności na okluzję.}
\label{tab:arch_map}
\end{table}

\subsection{Parametry funkcji straty}

\begin{table}[H]
\centering
\small
\begin{tabular}{lcc}
\toprule
\textbf{Komponenta straty} & \textbf{Model} & \textbf{Konfiguracja} \\
\midrule
Strata ArcMargin & Aligned\_Pretrained, CBAM\_Block & $m=0.5$, $s=30.0$ \\
Strata ArcMargin & Aligned\_Pretrained\_Transformers & $m=0.5$, $s=30.0$ (gałąź główna) \\
Strata Transformer (CE) & Aligned\_Pretrained\_Transformers & $\alpha=0.4$ w stracie złożonej \\
\bottomrule
\end{tabular}
\caption{Konfiguracja funkcji straty dla modeli obecnych w repozytorium.}
\label{tab:loss_params}
\end{table}


\section{Definicja metryk ewaluacji}
\label{app:metrics}

Dla zestawu testowego z $N$ całkowitymi zapytaniami sondy:

\begin{equation}
\text{Dokładność Rank-1} = \frac{1}{N} \sum_{i=1}^{N} \mathbb{1}[\text{top}\_1[i] = \text{gt}[i]]
\end{equation}

\begin{equation}
\text{Dokładność Rank-3} = \frac{1}{N} \sum_{i=1}^{N} \mathbb{1}[\text{gt}[i] \in \{\text{top}\_1, \text{top}\_2, \text{top}\_3\}[i]]
\end{equation}

gdzie $\mathbb{1}[\cdot]$ jest funkcją wskaźnika, a gt[i] jest rzeczywistą tożsamością dla zapytania $i$.

\section{Statystyki datasetu}
\label{app:dataset}

\begin{table}[H]
\centering
\small
\begin{tabular}{lrr}
\toprule
\textbf{Podział} & \textbf{Obrazy} & \textbf{Tożsamości} \\
\midrule
Trenowanie & 442 710 & 9 514 \\
Test (Galeria) & 24 227 & 1 057 \\
Test (Sonda) & 24 227 & 1 057 \\
\bottomrule
\end{tabular}
\caption{Statystyki podziału datasetu CASIA-WebFace.}
\label{tab:dataset_stats}
\end{table}


\end{document}